\scanpage{164}
\mspage{113}
% archive sheet number: 164
$U'_{/F'}$ est un objet tot. $W$-asphérique de $\mathrm{Cat}$).

\uline{Corollaire 3} \inserted{(du th.~3)} Si $i$\inserted{$:U\unsure{-|\to}X$} est une immersion ouverte ou fermée $W$-dewyérienne (resp. \unsure{totalement} $W$-dewyérienne), alors pour toute flèche $p:U-U'$ dans $\mathrm{Cat}$, $i':U'\to X'=U'\amalg_{\unsure{U|U'}}\unsure{X'|X}$ l'est aussi. [\deleted{\unsure{C'est bien}}\inserted{NB} le cas non respé est essentiellement le théorème 3. Il \unsure{ne|en} dépend \unsure{d'aucun} axiome supplémentaire sur $W$, mais n'a d'intérêt qu'en vue des \unsure{équivalences} \textcircled{a} et \textcircled{b} (p.~108)\inserted{ qui en résultent}, et qui sont valables moyennant $W(7\mathrm{bis})$ resp. $W(7)$ seulement. ]

\uline{Commentaire} de \uline{motivation} : les immersions ouvertes ou fermées $W$-dewyériennes, ou \deleted{\unsure{les}} tot\textsuperscript{t} $W$-dewyériennes, \unsure{étant|sont} \deleted{\unsure{par th. st.}} stables par cochangement de base (d'après le corollaire précédent), mais n'ont pas du tout l'air de l'être par composition. A cet égard, les \lq\lq immersions de Dwyer\rq\rq{} (où on suppose l'existence de l'adjoint à droite, \deleted{\unsure{resp.}} \deleted{\unsure{à gauche}}, de $\unsure{v_0:|v_0\in} U-\overline U$, dans le cas d'une imm. ouverte) se comportent mieux, étant stables par composition. (\uline{NB} Il ne semble pas
